\section{An\'alisis y correcci\'on de preguntas defectuosas}

\subsection{Criterios de auditor\'ia tras el ensayo}
Durante el ensayo de juego de roles se detectaron interrupciones en la fluidez de varias respuestas: el participante omiti\'o partes de las preguntas compuestas, respondi\'o con monos\'ilabos ante formulaciones cerradas o repiti\'o el riesgo sugerido por el entrevistador. Con base en las recomendaciones de elicitaci\'on de Pohl \cite{pohl2025}, el temario IREB \cite{ireb2026} y los criterios de no ambig\"uedad de ISO/IEC/IEEE 29148:2018 \cite{iso29148}, se auditaron las 15 preguntas originales y se identificaron \textbf{ocho preguntas defectuosas} clasificadas en cuatro categor\'ias: compuestas (dobles o triples), cerradas sin necesidad, inductivas y centradas en la soluci\'on.

\subsection{Cat\'alogo de preguntas que fallaron y su versi\'on corregida}

{\footnotesize
\noindent\begin{tabularx}{\linewidth}{|p{0.06\linewidth}|p{0.27\linewidth}|p{0.17\linewidth}|Y|}
\hline
\thdr{ID} & \thdr{Pregunta original (fall\'o en ensayo)} & \thdr{Defecto detectado} & \thdr{Versi\'on corregida (abierta, neutral y simple)}\\
\hline
\textbf{P3} & ?`En qu\'e horarios del d\'ia se concentran las mayores filas, cu\'antos clientes llegan a esperar y qu\'e pasos le toman m\'as tiempo en horas pico? & \textbf{Pregunta triple:} solicita horarios, volumen de espera y cuellos de botella en un solo enunciado. & \textbf{P3a:} ?`En qu\'e franjas horarias se acumula mayor cantidad de clientes en el mostrador? \newline \textbf{P3b:} Durante esas horas pico, ?`qu\'e pasos de la atenci\'on consumen m\'as tiempo?\\
\hline
\rowcolor{rowgray} \textbf{P4} & ?`C\'omo localiza un producto en el sistema y en bodega, y qu\'e tan frecuente es que el sistema indique existencias pero f\'isicamente est\'e agotado? & \textbf{Pregunta doble:} mezcla el mecanismo de b\'usqueda con la frecuencia de desfases de inventario. & \textbf{P4a:} ?`Qu\'e pasos sigue para localizar un producto en el sistema y ubicarlo f\'isicamente en bodega? \newline \textbf{P4b:} ?`Qu\'e ocurre cuando el registro indica disponibilidad pero el producto no aparece en percha?\\
\hline
\textbf{P5} & ?`Qu\'e herramienta usan hoy para registrar ventas y emitir facturas, qu\'e le genera mayor fatiga durante el cobro y qu\'e pas\'o la \'ultima vez que una venta se complic\'o? & \textbf{Pregunta triple:} combina herramientas actuales, fatiga operativa y relato de un incidente cr\'itico. & \textbf{P5a:} ?`Qu\'e herramientas utiliza actualmente para registrar ventas y emitir comprobantes? \newline \textbf{P5b:} Cu\'enteme qu\'e ocurri\'o la \'ultima vez que una venta se retras\'o o se complic\'o en caja.\\
\hline
\rowcolor{rowgray} \textbf{P8} & ?`Hay clientes que buscan siempre a un vendedor espec\'ifico para que los asesore y por qu\'e cree que ocurre eso? & \textbf{Cerrada al inicio y doble:} induce una respuesta de s\'i/no antes de explorar la causa. & ?`De qu\'e manera influye la experiencia de cada vendedor cuando un cliente solicita asesor\'ia t\'ecnica en el mostrador?\\
\hline
\textbf{P11} & Cuando el producto solicitado est\'a agotado, ?`en qu\'e criterios t\'ecnicos se basa para ofrecer un sustituto sin arriesgarse a provocar un cortocircuito o un reclamo? & \textbf{Inductiva:} condiciona la respuesta mencionando expl\'icitamente \enquote{cortocircuito} y \enquote{reclamo}. & Cuando un producto solicitado est\'a agotado, ?`qu\'e criterios eval\'ua usted para decidir qu\'e alternativa ofrecer al cliente?\\
\hline
\rowcolor{rowgray} \textbf{P12} & ?`Suele ofrecer productos complementarios (cinta aislante, cajas t\'ermicas, terminales) y c\'omo sabe cu\'ales sugerir en cada compra? & \textbf{Cerrada, doble y sesgada:} enumera ejemplos concretos que orientan al entrevistado. & ?`C\'omo decide qu\'e productos adicionales o complementarios sugerir junto con el material principal que pide el cliente?\\
\hline
\textbf{P14} & Si el sistema mostrara sugerencias de sustitutos o complementos al vender, ?`en qu\'e momento le resultar\'ia \'util verlas y qu\'e datos necesita en pantalla para confiar en esa recomendaci\'on? & \textbf{Centrada en la soluci\'on y doble:} asume una interfaz hipot\'etica antes de validar la necesidad real. & \textbf{P14a:} Al momento de atender en mostrador, ?`qu\'e informaci\'on t\'ecnica necesita verificar para estar seguro de que un reemplazo es compatible? \newline \textbf{P14b:} ?`En qu\'e instante de la venta le resulta necesario consultar esa informaci\'on?\\
\hline
\rowcolor{rowgray} \textbf{P15} & ?`Con qu\'e frecuencia falla el internet o se satura el servicio del SRI durante la semana y qu\'e procedimiento siguen en ese instante para no detener el cobro? & \textbf{Pregunta doble:} une frecuencia de fallos externos con protocolo de contingencia operativa. & \textbf{P15a:} ?`Con qu\'e frecuencia experimentan interrupciones de internet o lentitud en el SRI durante la semana? \newline \textbf{P15b:} Cuando ocurre una interrupci\'on de ese tipo, ?`qu\'e procedimiento siguen en caja para continuar atendiendo?\\
\hline
\end{tabularx}
}

\subsection{Ajustes aplicados para la entrevista real de campo (PE-U2 y PFC)}
La auditor\'ia demostr\'o que agrupar varias consultas en un solo enunciado provoca que el participante omita detalles relevantes sobre los tiempos de espera y los procedimientos de contingencia. Para las sesiones reales en los establecimientos de suministros el\'ectricos de Quevedo se aplicar\'an preguntas at\'omicas de una sola idea, se retirar\'an t\'erminos inductivos y se usar\'an repreguntas breves (\emph{\enquote{?`me puede dar un ejemplo reciente?}}) solo cuando se requiera ampliar una respuesta.
